\documentclass[12pt,letterpaper,oneside]{report}

% ================================================================
% Informe autocontenido del Laboratorio 01
% Compila con: pdflatex INFORME_LABORATORIO_01.tex (dos veces)
% No requiere imágenes externas: los diagramas se dibujan con LaTeX.
% ================================================================

\usepackage[utf8]{inputenc}
\usepackage[T1]{fontenc}
% `provide=*` permite usar la configuración moderna de babel aun cuando una
% instalación mínima de TeX no incluya el archivo heredado spanish.ldf.
\usepackage[provide=*,spanish]{babel}
\usepackage{lmodern}
\usepackage{geometry}
\usepackage{xcolor}
\usepackage{graphicx}
\usepackage{amsmath}
\usepackage{array}
\usepackage{booktabs}
\usepackage{tabularx}
\usepackage{longtable}
\usepackage{verbatim}
\usepackage{float}
\usepackage{url}
\usepackage{hyperref}

\geometry{top=2.3cm,bottom=2.4cm,left=2.5cm,right=2.5cm}
\setlength{\parindent}{0pt}
\setlength{\parskip}{0.65em}
\setlength{\emergencystretch}{3em}
\raggedbottom
\renewcommand{\arraystretch}{1.2}
\setcounter{secnumdepth}{3}
\setcounter{tocdepth}{2}

\definecolor{uvgblue}{HTML}{123B65}
\definecolor{compipurple}{HTML}{6D4AFF}
\definecolor{compicyan}{HTML}{007F8B}
\definecolor{softblue}{HTML}{EAF3FA}
\definecolor{softgray}{HTML}{F3F5F7}
\definecolor{darkgray}{HTML}{343A40}
\definecolor{okgreen}{HTML}{19733B}
\definecolor{errorred}{HTML}{A6232D}

\hypersetup{
  colorlinks=true,
  linkcolor=uvgblue,
  urlcolor=compicyan,
  citecolor=compipurple,
  pdftitle={Laboratorio 01 -- Analizador léxico y sintáctico de Compiscript},
  pdfauthor={DragonsSlayers 2.0},
  pdfsubject={Construcción de Compiladores CC-3032}
}

\newcommand{\proyecto}{Analizador de Compiscript}
\newcommand{\archivo}[1]{\texttt{#1}}
\newcommand{\regla}[1]{\textsf{#1}}
\newcommand{\ANTLR}{ANTLR 4}
\newcommand{\cps}{\texttt{.cps}}
\newcommand{\separador}{\par\noindent\color{uvgblue}\rule{\textwidth}{0.6pt}\color{black}\par}

% Cajas simples y portables, sin paquetes externos.
\newenvironment{notatecnica}
  {\begin{center}\begin{minipage}{0.92\textwidth}\color{darkgray}\hrule\vspace{0.5em}\textbf{Nota técnica.} }
  {\vspace{0.5em}\hrule\end{minipage}\end{center}}

\begin{document}
\hypersetup{pageanchor=false}

% ----------------------------------------------------------------
% Portada
% ----------------------------------------------------------------
\begin{titlepage}
  \centering
  \vspace*{1.0cm}
  {\Large\bfseries Universidad del Valle de Guatemala\par}
  \vspace{0.35cm}
  {\large Licenciatura en Ingeniería en Ciencias de la Computación\par}
  \vspace{0.2cm}
  {\large Construcción de Compiladores -- CC-3032, sección 10\par}

  \vfill

  \color{uvgblue}\rule{0.88\textwidth}{1.2pt}\color{black}
  \vspace{0.8cm}

  {\Huge\bfseries Laboratorio 01\par}
  \vspace{0.35cm}
  {\LARGE Uso de generadores de analizadores léxicos y sintácticos\par}
  \vspace{0.7cm}
  {\Huge\color{compipurple}\bfseries \proyecto\par}
  \vspace{0.3cm}
  {\Large Documentación teórica y práctica\par}

  \vspace{0.8cm}
  \color{uvgblue}\rule{0.88\textwidth}{1.2pt}\color{black}

  \vfill

  \begin{tabular}{>{\bfseries}rl}
    Catedrático: & Ingeniero Carlos Valdez \\
    Grupo: & DragonsSlayers 2.0 \\
    Integrantes: & Pablo Daniel Barillas Moreno -- 22193 \\
                 & Hugo Daniel Barillas Ajín -- 23556 \\
                 & Ernesto Ascencio -- 23009 \\
  \end{tabular}

  \vfill
  {\large Guatemala, 2026\par}
\end{titlepage}

\hypersetup{pageanchor=true}
\pagenumbering{roman}

\chapter*{Resumen}
\addcontentsline{toc}{chapter}{Resumen}

Este documento presenta la fundamentación, diseño, implementación y validación del
Laboratorio 01 de Construcción de Compiladores. El producto desarrollado es una
herramienta que analiza archivos de Compiscript con extensión \cps. La solución
utiliza una gramática combinada de \ANTLR{} y el generador \archivo{antlr4ts-cli}
para producir un lexer y un parser TypeScript reales. El mismo núcleo se consume
desde una interfaz web construida con React y Vite, una aplicación de consola y una
aplicación de escritorio empaquetada con Electron.

La herramienta fuerza la tokenización completa, conserva los tokens reconocibles
después de errores léxicos, ejecuta el parser desde la regla \regla{program} y aplica
\archivo{DefaultErrorStrategy} para recuperarse de errores sintácticos. Los
diagnósticos se normalizan al español e incluyen fase, línea, columna y símbolo
relacionado. También se agrupan caracteres inválidos contiguos, se eliminan
duplicados, se filtran cascadas derivadas y se limita la salida a 100 diagnósticos
por fase sin detener el procesamiento.

El resultado incluye tabla de tokens, errores separados por fase, resumen de
aceptación y árbol de parseo textual y visual. La suite automatizada contiene 24
pruebas que cubren construcciones válidas, recuperación ante múltiples errores,
cadenas incorrectas, sincronización de ejemplos, límites de diagnósticos y los ocho
casos de complejidad baja y media definidos por la rúbrica. El alcance se limita
deliberadamente a análisis léxico y sintáctico: no hay análisis semántico,
interpretación ni generación de código.

\textbf{Palabras clave:} compiladores, análisis léxico, análisis sintáctico, ANTLR,
Compiscript, TypeScript, recuperación de errores.

\tableofcontents
\listoffigures
\listoftables
\clearpage
\pagenumbering{arabic}

% ================================================================
\chapter{Introducción}
% ================================================================

Un compilador o intérprete necesita transformar texto plano en una representación
estructurada antes de poder comprobar significados o ejecutar instrucciones. Las dos
primeras etapas de esta transformación son el análisis léxico y el análisis
sintáctico. El primero agrupa caracteres en unidades llamadas \emph{tokens}; el
segundo determina si la secuencia de tokens puede derivarse de una gramática formal.

El laboratorio estudia estas etapas mediante una herramienta generadora. En lugar de
programar manualmente una máquina de estados y un parser, el lenguaje se describe en
un archivo de gramática: \archivo{Compiscript.g4}. ANTLR genera las clases que
reconocen la entrada y la aplicación se ocupa de coordinarlas, recolectar resultados,
mejorar los mensajes y presentarlos de forma útil.

Compiscript es un lenguaje educativo inspirado en TypeScript. Incluye declaraciones,
tipos básicos, arreglos, expresiones con precedencia, funciones, clases y estructuras
de control. La aplicación permite cargar código propio o usar tres casos demostrativos
integrados. Adicionalmente, \path{examples/rubric} aporta ocho entradas reproducibles
que cubren los cuatro escenarios de complejidad baja y los cuatro de complejidad
media solicitados en la evaluación.

\section{Objetivo general}

Diseñar e implementar un analizador léxico y sintáctico para archivos Compiscript
mediante un generador de analizadores, con una interfaz que muestre tokens,
diagnósticos comprensibles y el árbol de parseo sin detenerse en el primer error
recuperable.

\vspace{8cm}

\section{Objetivos específicos}

\begin{itemize}
  \item Definir la sintaxis concreta de Compiscript en una gramática ANTLR.
  \item Generar e integrar un lexer y un parser TypeScript a partir de esa gramática.
  \item Mostrar tokens con tipo, lexema, posición y canal.
  \item Distinguir errores léxicos de errores sintácticos y traducir los mensajes
        técnicos más comunes a español.
  \item Continuar el análisis ante entradas incorrectas y controlar las cascadas de
        diagnósticos.
  \item Visualizar y exportar los resultados desde una interfaz gráfica.
  \item Ofrecer ejecución alternativa mediante CLI y empaquetado de escritorio.
  \item Verificar el comportamiento mediante pruebas automatizadas y casos \cps.
\end{itemize}

\section{Alcance y exclusiones}

El proyecto reconoce la estructura léxica y sintáctica del lenguaje. Un archivo se
considera aceptado cuando las listas de errores léxicos y sintácticos están vacías.
El árbol producido demuestra qué estructura reconoció el parser.

No se construye una tabla de símbolos, no se resuelven nombres, no se comprueban
tipos, no se evalúan expresiones y no se genera código intermedio, máquina virtual o
ensamblador. Por ejemplo, una asignación entre tipos incompatibles puede ser
sintácticamente válida aunque sería rechazada en una futura fase semántica.

% ================================================================
\chapter{Fundamentos teóricos}
% ================================================================

\section{De caracteres a estructura}

El frente de un compilador puede representarse como una composición de funciones:

\[
  \text{caracteres}\xrightarrow{\text{lexer}}\text{tokens}
  \xrightarrow{\text{parser}}\text{árbol sintáctico}.
\]

Esta separación reduce complejidad. El parser no necesita distinguir cada carácter
de una palabra reservada; recibe un token como \archivo{IF} o \archivo{Identifier}.
El lexer, a su vez, no decide si un \archivo{IF} aparece en una posición gramatical
válida.

\section{Análisis léxico}

Un \textbf{lexema} es el fragmento concreto de entrada, por ejemplo
\archivo{contador}; un \textbf{token} es su categoría, en este caso
\archivo{Identifier}. Un patrón describe el conjunto de lexemas de una categoría.
En Compiscript, un identificador comienza con letra o guion bajo y continúa con
letras, dígitos o guiones bajos.

El lexer generado aplica prioridades y la política de coincidencia más larga. Por
ello, las palabras reservadas se declaran antes de \archivo{Identifier}, y los
operadores de dos caracteres se declaran antes que sus versiones de uno. Así,
\archivo{<=} se reconoce como una unidad y no como \archivo{<} seguida de
\archivo{=}.

Los espacios, saltos de línea y comentarios se envían a reglas con acción
\archivo{skip}; no llegan al parser. Sin embargo, el lexer actualiza línea y columna
para que los tokens y errores siguientes mantengan una ubicación correcta.

\section{Gramáticas libres de contexto}

Una gramática libre de contexto se describe como una cuádrupla
$G=(V,\Sigma,P,S)$, donde $V$ es el conjunto de no terminales, $\Sigma$ el de
terminales, $P$ las producciones y $S$ el símbolo inicial. En el proyecto,
\regla{program} es el símbolo inicial y exige cero o más sentencias seguidas de EOF:

\begin{center}
\fbox{\texttt{program : statement* EOF ;}}
\end{center}

Los terminales son los tokens producidos por el lexer; los no terminales son reglas
como \regla{statement}, \regla{functionDeclaration} o \regla{expression}. El parser
construye un árbol cuyos nodos internos corresponden a reglas y cuyas hojas
corresponden a tokens.

\section{Precedencia y asociatividad}

Una gramática de expresiones debe evitar que \archivo{1 + 2 * 3} se interprete como
\archivo{(1 + 2) * 3}. Compiscript separa cada nivel de precedencia en una regla.
Cuanto más cerca está una regla de \regla{primaryExpression}, mayor es su prioridad.

\begin{table}[H]
\centering
\caption{Precedencia de expresiones, de menor a mayor.}
\label{tab:precedencia}
\begin{tabularx}{\textwidth}{>{\bfseries}c l X}
\toprule
Nivel & Regla & Operadores o forma \\
\midrule
1 & Asignación & \archivo{=} \\
2 & Condicional & \archivo{? :} \\
3 & OR lógico & \archivo{||} \\
4 & AND lógico & \archivo{\&\&} \\
5 & Igualdad & \archivo{==}, \archivo{!=} \\
6 & Relacional & \archivo{<}, \archivo{<=}, \archivo{>}, \archivo{>=} \\
7 & Aditiva & \archivo{+}, \archivo{-} \\
8 & Multiplicativa & \archivo{*}, \archivo{/}, \archivo{\%} \\
9 & Unaria & \archivo{!}, \archivo{-}, \archivo{+} \\
10 & Primaria & literales, identificadores, llamadas, índices y propiedades \\
\bottomrule
\end{tabularx}
\end{table}

\vspace{10cm}

\section{ANTLR 4 y Adaptive LL(*)}

ANTLR recibe una gramática declarativa y genera reconocedores para un lenguaje
objetivo. En este laboratorio, \archivo{antlr4ts-cli} genera TypeScript. El parser de
ANTLR 4 utiliza predicción \textbf{Adaptive LL(*)}: analiza la entrada de izquierda a
derecha, construye una derivación izquierda y adapta su predicción con autómatas
internos. Puede examinar la cantidad de contexto necesaria para distinguir
alternativas.

No debe confundirse con un parser SLR. La implementación no posee tablas
\archivo{ACTION/GOTO} ni registra operaciones \emph{shift/reduce}. Su evidencia de
ejecución es la invocación a las clases generadas y la construcción del árbol ANTLR.

\section{Recuperación de errores}

Detenerse en el primer error dificulta corregir un archivo grande. Una estrategia de
recuperación intenta alcanzar un punto estable y continuar. La estrategia
\archivo{DefaultErrorStrategy} de ANTLR puede:

\begin{itemize}
  \item insertar conceptualmente un token faltante;
  \item eliminar un token sobrante;
  \item sincronizar el flujo con tokens esperados por reglas posteriores.
\end{itemize}

La recuperación no convierte una entrada inválida en válida: permite construir un
árbol parcial y recolectar más diagnósticos. Es necesario controlar cascadas, porque
un error inicial puede provocar mensajes secundarios que no representan problemas
independientes.

\subsection{Propagación y cascadas de errores}

Cuando el lexer descarta un carácter no reconocido, el parser recibe un flujo de
tokens distinto del texto original. Esa ausencia puede provocar un diagnóstico
sintáctico secundario aunque el usuario haya cometido un solo error léxico. De forma
similar, durante la recuperación sintáctica, la inserción o eliminación conceptual de
un token puede modificar la interpretación de construcciones posteriores. A este
fenómeno se le denomina \emph{propagación} o \emph{cascada de errores}.

La implementación no crea errores artificiales: conserva los diagnósticos producidos
por ANTLR y aplica medidas para presentar únicamente información útil. Agrupa
caracteres inválidos contiguos, elimina duplicados exactos, evita volver a reportar
símbolos ya cubiertos por una cadena inválida y filtra errores del parser ubicados
inmediatamente después de un fragmento reportado por el lexer. Las cascadas
sintácticas lejanas no pueden clasificarse causalmente con certeza, por lo que se
conservan hasta un máximo de cien diagnósticos por fase.

% ================================================================
\chapter{Especificación de Compiscript}
% ================================================================

\section{Unidades léxicas}

\begin{table}[H]
\centering
\caption{Familias principales de tokens.}
\label{tab:tokens}
\begin{tabularx}{\textwidth}{>{\bfseries}l X}
\toprule
Familia & Elementos representativos \\
\midrule
Declaración & \archivo{let}, \archivo{var}, \archivo{const}, \archivo{function}, \archivo{class} \\
Control & \archivo{if}, \archivo{else}, \archivo{while}, \archivo{for}, \archivo{foreach}, \archivo{switch}, \archivo{try}, \archivo{catch} \\
Tipos & \archivo{integer}, \archivo{string}, \archivo{boolean} e identificadores de clase \\
Literales & enteros, cadenas, \archivo{true}, \archivo{false}, \archivo{null} y arreglos \\
Operadores & aritméticos, relacionales, igualdad, lógicos, asignación y ternario \\
Delimitadores & paréntesis, llaves, corchetes, coma, punto, dos puntos y punto y coma \\
Ignorados & espacios, tabulaciones, saltos de línea y comentarios de línea o bloque \\
\bottomrule
\end{tabularx}
\end{table}

Las cadenas admiten los escapes \archivo{\textbackslash"},
\archivo{\textbackslash\textbackslash}, \archivo{\textbackslash n},
\archivo{\textbackslash r} y \archivo{\textbackslash t}. Una comilla inicial sin
cierre en la misma línea o una secuencia no permitida produce un diagnóstico léxico
especializado.

\section{Declaraciones y tipos}

\archivo{let} y \archivo{var} permiten anotación de tipo e inicializador opcionales.
\archivo{const} requiere inicializador. Los tipos básicos son \archivo{integer},
\archivo{string} y \archivo{boolean}; un identificador puede representar una clase.
Cada sufijo \archivo{[]} agrega una dimensión de arreglo.

\begin{verbatim}
let contador: integer = 0;
var mensaje: string;
const ACTIVO: boolean = true;
let matriz: integer[][] = [[1, 2], [3, 4]];
\end{verbatim}

\section{Funciones y clases}

Una función puede declarar parámetros tipados y un tipo de retorno. La gramática
acepta recursión y funciones anidadas porque una declaración de función es una
sentencia válida dentro de un bloque. Las clases admiten miembros, herencia mediante
dos puntos, acceso con \archivo{this} e instanciación con \archivo{new}.

\begin{verbatim}
function factorial(n: integer): integer {
  if (n <= 1) return 1;
  return n * factorial(n - 1);
}

class Perro : Animal {
  function hablar(): string {
    return this.nombre + " ladra.";
  }
}
\end{verbatim}

\section{Control de flujo}

El lenguaje cubre bloques y ciclos \archivo{while}, \archivo{do-while},
\archivo{for} y \archivo{foreach}. La selección emplea \archivo{if/else} o
\archivo{switch}, \archivo{case} y \archivo{default}; el manejo sintáctico usa
\archivo{try/catch}. También existen \archivo{break}, \archivo{continue} y
\archivo{return}. La palabra
``sintáctico'' es importante: el parser acepta la forma, pero no verifica, por
ejemplo, que \archivo{break} aparezca semánticamente dentro de un ciclo.

\section{Accesos y expresiones}

La regla del lado izquierdo parte de un átomo primario y admite una sucesión de
sufijos: llamada \archivo{(...)}, índice \archivo{[...]}, o propiedad
\archivo{.identificador}. Esto permite reconocer construcciones como
\archivo{obj.metodo(lista[0]).propiedad}. El agrupamiento con paréntesis y los niveles
de la Tabla~\ref{tab:precedencia} preservan la estructura esperada.

% ================================================================
\chapter{Arquitectura y diseño}
% ================================================================

\section{Arquitectura general}

La lógica del analizador se concentra en \archivo{src/lib}; no está duplicada entre
las interfaces. React y la CLI llaman a \archivo{analyzeInput}. Electron aloja la
aplicación web ya compilada, por lo que también reutiliza el mismo núcleo.

\begin{figure}[H]
\centering
\setlength{\unitlength}{1mm}
\begin{picture}(150,70)
  \put(2,50){\framebox(34,12){\shortstack{Archivo \.cps}}}
  \put(43,50){\framebox(42,12){\shortstack{Núcleo ANTLR\\\archivo{src/lib}}}}
  \put(92,50){\framebox(55,12){\shortstack{Resultado normalizado\\\archivo{AnalyzeResult}}}}
  \put(36,56){\vector(1,0){7}}
  \put(85,56){\vector(1,0){7}}

  \put(4,10){\framebox(37,15){\shortstack{React + Vite\\interfaz web}}}
  \put(56,10){\framebox(37,15){\shortstack{CLI TypeScript\\terminal}}}
  \put(108,10){\framebox(37,15){\shortstack{Electron\\Windows x64}}}

  \put(111,50){\vector(-3,-2){60}}
  \put(119,50){\vector(-1,-1){25}}
  \put(128,50){\vector(1,-3){8}}
\end{picture}
\caption{Arquitectura de reutilización del núcleo de análisis.}
\label{fig:arquitectura}
\end{figure}

\vspace{8cm}

\section{Generación e integración}

La gramática es la fuente de verdad del lenguaje. El comando
\archivo{npm run generate} ejecuta \archivo{antlr4ts} con opciones de visitor y
listener, y escribe los artefactos en \archivo{src/generated}. Estos archivos se
incluyen en el repositorio para que Java sea necesario únicamente al regenerarlos.

\begin{figure}[H]
\centering
\setlength{\unitlength}{1mm}
\begin{picture}(150,42)
  \put(1,16){\framebox(36,13){\shortstack{\archivo{Compiscript.g4}\\fuente}}}
  \put(43,16){\framebox(30,13){\shortstack{\archivo{antlr4ts-cli}\\Java}}}
  \put(79,16){\framebox(42,13){\shortstack{Lexer, parser,\\listener y visitor}}}
  \put(127,16){\framebox(22,13){\shortstack{TypeScript\\compilado}}}
  \put(37,22){\vector(1,0){6}}
  \put(73,22){\vector(1,0){6}}
  \put(121,22){\vector(1,0){6}}
\end{picture}
\caption{Flujo de generación de analizadores.}
\label{fig:generacion}
\end{figure}

Los archivos generados no deben modificarse manualmente: una regeneración los
sobrescribe. Para cambiar el lenguaje se edita primero la gramática, luego se genera
de nuevo y finalmente se ejecutan las verificaciones.

\section{Flujo en tiempo de ejecución}

\begin{figure}[H]
\centering
\setlength{\unitlength}{1mm}
\begin{picture}(150,104)
  \put(45,90){\framebox(60,10){Texto Compiscript}}
  \put(45,72){\framebox(60,10){\archivo{ANTLRInputStream}}}
  \put(45,54){\framebox(60,10){\archivo{CompiscriptLexer}}}
  \put(45,36){\framebox(60,10){\archivo{CommonTokenStream.fill()}}}
  \put(3,13){\framebox(51,12){\shortstack{Modo lexer\\tokens + errores léxicos}}}
  \put(96,13){\framebox(51,12){\shortstack{Modo parser: \archivo{program()}\\errores + árbol}}}
  \put(75,90){\vector(0,-1){8}}
  \put(75,72){\vector(0,-1){8}}
  \put(75,54){\vector(0,-1){8}}
  \put(60,36){\vector(-2,-1){25}}
  \put(90,36){\vector(2,-1){25}}
  \put(109,28){\makebox(33,5){\scriptsize \archivo{seek(0)}}}
\end{picture}
\caption{Pipeline de análisis y bifurcación por modo.}
\label{fig:pipeline}
\end{figure}

Los pasos concretos son:

\begin{enumerate}
  \item Crear un \archivo{ANTLRInputStream} a partir del texto.
  \item Instanciar \archivo{CompiscriptLexer} y sustituir sus listeners por uno
        recolector.
  \item Crear \archivo{CommonTokenStream} y llamar \archivo{fill()} para forzar el
        recorrido completo.
  \item Convertir cada token, excepto EOF, a \archivo{TokenInfo}.
  \item Si el modo es léxico, terminar sin crear un parser.
  \item En modo completo, rebobinar con \archivo{seek(0)} y crear el parser.
  \item Configurar \archivo{DefaultErrorStrategy} e invocar \regla{program}.
  \item Convertir el árbol ANTLR a texto Lisp, texto indentado y nodos visuales.
  \item Filtrar diagnósticos derivados y calcular la aceptación y el resumen.
\end{enumerate}

\section{Modelo de resultado}

\archivo{AnalyzeResult} actúa como contrato entre el núcleo y sus consumidores.
Contiene el lenguaje, modo, aceptación, tokens, errores por fase, tres formas del
árbol, explicación y conteos. Cada token conserva índice, ID numérico, nombre
simbólico, texto, línea, columna base 1 y canal. Cada error conserva origen, posición,
mensaje, símbolo opcional y severidad.

Esta estructura desacopla el motor de la presentación. La UI puede renderizar tablas
y árboles, la CLI puede imprimir texto y las funciones de descarga pueden serializar
JSON o CSV sin volver a ejecutar el parser.

% ================================================================
\chapter{Implementación práctica}
% ================================================================

\section{Organización del repositorio}

\begin{verbatim}
.
|-- docs/                 Documentación y matriz de cumplimiento
|-- electron/main.cjs     Proceso principal de Electron
|-- examples/
|   |-- compiscript/      Casos integrados de la interfaz
|   |-- rubric/low/       Cuatro casos de complejidad baja
|   |-- rubric/medium/    Cuatro casos de complejidad media
|   `-- error_distribution/
|       |-- lexer_errors.cps
|       |-- parser_errors.cps
|       `-- lexer_parser_alternating_errors.cps
|-- src/
|   |-- __tests__/        Pruebas Vitest
|   |-- cli/run.ts        Entrada de consola
|   |-- generated/        Código producido por ANTLR
|   |-- grammars/         Compiscript.g4
|   |-- lib/              Análisis, errores, árboles y descargas
|   |-- ui/               App y componentes React
|   |-- main.tsx          Punto de entrada web
|   `-- styles.css        Diseño adaptable
|-- package.json          Scripts y metadatos de empaquetado
|-- tsconfig.json         TypeScript estricto
`-- vite.config.ts        Vite, React y polyfills
\end{verbatim}

\vspace{5cm}

\section{Gramática combinada}

\archivo{Compiscript.g4} contiene reglas sintácticas y léxicas. La regla superior
admite una lista de sentencias y exige llegar al final del archivo. Una versión
reducida de la estructura es:

\begin{verbatim}
grammar Compiscript;

program
  : statement* EOF
  ;

statement
  : variableDeclaration
  | constantDeclaration
  | functionDeclaration
  | classDeclaration
  | printStatement
  | block
  | ifStatement
  | whileStatement
  | doWhileStatement
  | forStatement
  | foreachStatement
  | tryCatchStatement
  | switchStatement
  | breakStatement
  | continueStatement
  | returnStatement
  | expressionStatement
  ;
\end{verbatim}

La expresión se divide por niveles en vez de usar un parser manual. La asignación
permite un lado izquierdo seguido de \archivo{=} y otra asignación, o delega a la
expresión condicional. La recursión a la derecha modela la asociatividad esperada de
la asignación.

\vspace{5cm}

\section{Integración en \archivo{analyze.ts}}

La función pública acepta el texto y un modo opcional:

\begin{verbatim}
export function analyzeInput(
  input: string,
  mode: AnalyzerMode = "parser"
): AnalyzeResult
\end{verbatim}

El uso de \archivo{tokenStream.fill()} es una decisión central: fuerza al lexer a
recorrer toda la entrada antes del parseo y hace posible capturar múltiples errores
léxicos. Después se recolecta la información de los tokens, se descarta EOF y se
obtiene cada nombre desde el vocabulario generado. En modo completo se ejecuta:

\begin{verbatim}
tokenStream.seek(0);
const parser = new CompiscriptParser(tokenStream);
parser.errorHandler = new DefaultErrorStrategy();
parser.buildParseTree = true;
const tree = parser.program();
\end{verbatim}

La decisión de aceptación es explícita:

\[
  accepted \iff |E_{lex}|=0 \land |E_{syn}|=0.
\]

Por tanto, producir tokens o un árbol parcial no implica que el archivo haya sido
aceptado.

\vspace{10cm}

\section{Recolección y normalización de errores}

\archivo{CollectingErrorListener} implementa \archivo{ANTLRErrorListener} tanto para
el lexer como para el parser. Convierte la columna interna base 0 a base 1 y genera un
objeto uniforme. Los mensajes comunes se transforman como muestra la tabla.

\begin{table}[H]
\centering
\caption{Normalización de mensajes de ANTLR.}
\label{tab:errores}
\begin{tabularx}{\textwidth}{>{\ttfamily}l X}
\toprule
Patrón interno & Presentación al usuario \\
\midrule
token recognition error & Carácter o lexema no reconocido \\
missing X at Y & Falta X antes de Y \\
extraneous input & Se encontró un token, pero se esperaba otro \\
mismatched input & El token no coincide con el conjunto esperado \\
no viable alternative & No se reconoció una construcción válida cercana \\
failed predicate & La construcción no cumple una regla de la gramática \\
\bottomrule
\end{tabularx}
\end{table}

El tratamiento adicional incluye:

\begin{itemize}
  \item \textbf{Cadenas:} si una comilla no tiene cierre no escapado en la línea, se
        reporta una cadena sin cerrar; si sí existe cierre, el fragmento completo se
        reporta como cadena inválida por sus escapes.
  \item \textbf{Agrupación:} errores de reconocimiento contiguos en la misma línea se
        unen; \archivo{@@@} produce un diagnóstico y no tres.
  \item \textbf{Deduplicación:} no se agrega dos veces la misma combinación de línea,
        columna, símbolo y mensaje.
  \item \textbf{Cobertura:} no se reporta de nuevo un carácter que ya está dentro del
        fragmento de una cadena inválida.
  \item \textbf{Filtrado cruzado:} se omiten errores sintácticos inmediatamente
        cercanos a un fragmento ya descartado por el lexer.
  \item \textbf{Límite:} el centésimo elemento se convierte en un aviso de límite;
        el análisis continúa aunque se omitan mensajes posteriores.
\end{itemize}

\begin{figure}[H]
\centering
\setlength{\unitlength}{1mm}
\begin{picture}(150,67)
  \put(2,48){\framebox(34,12){Error ANTLR}}
  \put(43,48){\framebox(44,12){\shortstack{Traducir y ubicar\\línea/columna}}}
  \put(94,48){\framebox(53,12){\shortstack{Especializar cadenas\\y fragmentos}}}
  \put(36,54){\vector(1,0){7}}
  \put(87,54){\vector(1,0){7}}
  \put(25,14){\framebox(44,12){\shortstack{Deduplicar y\\agrupar}}}
  \put(81,14){\framebox(44,12){\shortstack{Filtrar cascadas\\y limitar salida}}}
  \put(111,48){\vector(-2,-1){54}}
  \put(69,20){\vector(1,0){12}}
  \put(125,20){\vector(1,0){15}}
  \put(132,14){\framebox(17,12){Mostrar}}
\end{picture}
\caption{Proceso de normalización de diagnósticos.}
\label{fig:errores}
\end{figure}

\section{Representación del árbol}

El árbol se conserva de tres maneras:

\begin{enumerate}
  \item \archivo{parseTreeText}: salida Lisp de \archivo{toStringTree};
  \item \archivo{formattedParseTree}: recorrido indentado para descarga y copia;
  \item \archivo{parseTreeNodes}: nodos recursivos para la vista interactiva.
\end{enumerate}

La conversión visual recorre directamente \archivo{ParseTree}. Si el nodo es un
\archivo{ParserRuleContext}, usa el nombre de regla; si es terminal, usa su texto.
Esta ruta evita ambigüedades que aparecerían al intentar volver a interpretar la
cadena Lisp cuando un token literal es un paréntesis.

\begin{figure}[H]
\centering
\setlength{\unitlength}{1mm}
\begin{picture}(130,60)
  \put(50,48){\framebox(30,9){program}}
  \put(13,27){\framebox(34,9){statement}}
  \put(83,27){\framebox(34,9){EOF}}
  \put(65,48){\line(-2,-1){31}}
  \put(65,48){\line(2,-1){35}}
  \put(2,7){\framebox(26,9){LET}}
  \put(31,7){\framebox(29,9){Identifier}}
  \put(63,7){\framebox(22,9){SEMI}}
  \put(30,27){\line(-1,-1){15}}
  \put(30,27){\line(1,-1){15}}
  \put(30,27){\line(3,-1){44}}
\end{picture}
\caption{Ejemplo simplificado de árbol de parseo.}
\label{fig:arbol}
\end{figure}

\section{Interfaz web}

\archivo{App.tsx} coordina estado, caso seleccionado, texto de entrada, pestaña,
resultado y errores inesperados. La interfaz posee tres vistas:

\begin{itemize}
  \item \textbf{Análisis completo:} lexer, parser, diagnósticos y árbol;
  \item \textbf{Solo lexer:} tokens y errores léxicos sin instanciar el parser;
  \item \textbf{Documentación:} cobertura, comandos y alcance.
\end{itemize}

\archivo{InputEditor} restringe el selector a \cps, valida nuevamente la extensión,
lee el archivo como texto y permite editarlo. Los paneles presentan el estado,
conteos, tabla de tokens, grupos de errores y árbol colapsable. La hoja de estilos
utiliza rejillas y puntos de quiebre en 1100, 700 y 500 píxeles; los bloques de código,
tablas y árbol poseen desplazamiento o ajuste de palabra para evitar desbordamientos.

Las descargas disponibles son gramática \archivo{.g4}, entrada \cps, tokens CSV,
tokens JSON, árbol TXT en modo parser y resultado JSON. Se construyen con
\archivo{Blob} y una URL temporal del navegador.

\section{Compatibilidad con navegador}

\archivo{antlr4ts} depende indirectamente de \archivo{assert}, \archivo{util} y la
variable global \archivo{process}. Vite configura alias para los dos módulos y
\archivo{browserProcess.ts} instala el polyfill global antes de cargar el analizador.
Esta inicialización evita una pantalla en blanco por referencias de Node ausentes en
el navegador.

\section{Interfaz de consola}

\archivo{src/cli/run.ts} valida que exista un argumento, que el archivo tenga
extensión \cps y que esté disponible. Lee UTF-8, llama al mismo
\archivo{analyzeInput}, imprime conteos y errores con color, e informa si el análisis
continuó después de errores recuperables.

\begin{table}[H]
\centering
\caption{Códigos de salida de la CLI.}
\label{tab:exit}
\begin{tabular}{c l}
\toprule
Código & Significado \\
\midrule
0 & Archivo aceptado sin errores léxicos ni sintácticos \\
1 & Archivo rechazado o argumentos de usuario inválidos \\
2 & Excepción inesperada durante el análisis \\
\bottomrule
\end{tabular}
\end{table}

\section{Aplicación de escritorio}

Electron crea una ventana de $1440\times900$ píxeles, con mínimos de
$1000\times700$, y carga \archivo{dist/index.html}. La configuración de seguridad
desactiva la integración de Node, activa el aislamiento de contexto y mantiene el
\emph{sandbox}. El menú nativo se oculta.

\archivo{electron-builder} produce Windows x64 en dos modalidades:

\begin{center}
\archivo{Compiscript-Portable-<version>.exe}\qquad
\archivo{Compiscript-Setup-<version>.exe}
\end{center}

La aplicación final no requiere Node.js, npm ni Java en la computadora del usuario.
Java solo se requiere al regenerar la gramática durante el desarrollo.

% ================================================================
\chapter{Pruebas y validación}
% ================================================================

\section{Estrategia de pruebas}

La suite utiliza Vitest y contiene 24 casos: 13 instancias en el archivo principal
del analizador, 3 comprobaciones parametrizadas de sincronización entre los ejemplos
de disco y los textos mostrados en la interfaz, y 8 comprobaciones parametrizadas de
la matriz oficial de la rúbrica. Las pruebas no se limitan a verificar un valor
booleano; también inspeccionan tokens posteriores, mensajes, conteos, estructuras
exigidas y árbol.

\begingroup
\small
\renewcommand{\arraystretch}{1.1}
\begin{longtable}{p{0.29\textwidth} p{0.64\textwidth}}
\caption{Cobertura automatizada del proyecto.}\label{tab:pruebas}\\
\toprule
Área & Evidencia verificada \\
\midrule
\endfirsthead
\toprule
Área & Evidencia verificada \\
\midrule
\endhead
Programa integral & Aceptación, cero errores, tokens de clases y foreach, y regla de función en el árbol. \\
Cobertura válida & Tipos, matrices, expresiones, ciclos, switch, try/catch, clases, herencia, new, closures y recursión. \\
Errores léxicos & Dos caracteres inválidos, continuidad hasta tokens posteriores, mensajes normalizados y ausencia de parser en modo lexer. \\
Agrupación & \archivo{@@@} genera un solo error léxico y no una cascada sintáctica inmediata. \\
Cadenas & Diagnóstico específico para cadena sin cerrar y para escape inválido. \\
Límite léxico & 120 fragmentos producen 100 entradas visibles y todavía se reconoce el identificador final. \\
Recuperación sintáctica & Se reporta más de un error, no se exponen frases internas en inglés y se conserva el árbol. \\
Límite sintáctico & Se limita la lista y el árbol alcanza EOF. \\
Sincronización & Los tres archivos \cps coinciden con los casos embebidos en la interfaz. \\
Matriz de la rúbrica & Ocho archivos comprueban complejidad baja/media y los mínimos de errores por fase. \\
\bottomrule
\end{longtable}
\endgroup

\section{Casos demostrativos}

El caso válido combina constantes, recursión, clases, herencia, arreglos,
\archivo{foreach} y \archivo{try/catch}. Debe producir cero errores y un árbol desde
\regla{program} hasta EOF.

El caso léxico contiene dos símbolos no definidos y una sentencia posterior:

\begin{verbatim}
let primero: integer = 10 @ 2;
let segundo: integer = 20 # 4;
print(primero);
\end{verbatim}

La prueba espera dos errores y confirma que el token \archivo{print} todavía existe,
lo cual demuestra continuidad del escaneo.

El caso sintáctico omite puntos y coma y un paréntesis de cierre:

\begin{verbatim}
let nombre: string = "Compiscript"
const total: integer = 10

if (total > 5 {
  print(nombre)
}
\end{verbatim}

El parser debe recuperar más de un error y conservar una representación del árbol.

\section{Matriz de casos exigidos por la rúbrica}

La carpeta \path{examples/rubric} transforma cada criterio de evaluación en una
entrada independiente. Los cuatro casos bajos contienen variables de tipo
\archivo{integer}, \archivo{string} y \archivo{boolean}; una constante; suma y
multiplicación; un condicional; un ciclo \archivo{while}; y un \archivo{foreach}.
Los cuatro casos medios conservan esos elementos y agregan un arreglo declarado y
utilizado, dos clases, dos objetos instanciados, dos funciones declaradas y dos
llamadas.

\begin{table}[H]
\centering
\caption{Resultados verificados de la matriz de la rúbrica.}
\label{tab:matriz-rubrica}
\begin{tabularx}{\textwidth}{>{\ttfamily}X r r r r}
\toprule
\normalfont Caso & Tokens & Léxicos & Sintácticos & Salida \\
\midrule
low/valid.cps & 112 & 0 & 0 & 0 \\
low/lexer\_errors.cps & 112 & 3 & 0 & 1 \\
low/parser\_errors.cps & 109 & 0 & 3 & 1 \\
low/lexer\_parser\_errors.cps & 110 & 2 & 2 & 1 \\
medium/valid.cps & 263 & 0 & 0 & 0 \\
medium/lexer\_errors.cps & 263 & 3 & 0 & 1 \\
medium/parser\_errors.cps & 260 & 0 & 3 & 1 \\
medium/lexer\_parser\_errors.cps & 261 & 2 & 2 & 1 \\
\bottomrule
\end{tabularx}
\end{table}

Los códigos de salida cero corresponden a aceptación sin errores. El código uno en
los otros seis archivos confirma el rechazo intencional. Los casos exclusivamente
léxicos no originan errores sintácticos derivados; los exclusivamente sintácticos no
contienen símbolos fuera del vocabulario; y los mixtos alcanzan exactamente dos
diagnósticos de cada fase. Todos presentan fase, línea, columna, símbolo o token y
mensaje en español.

La automatización en \archivo{rubric.examples.test.ts} lee los archivos reales,
ejecuta \archivo{analyzeInput}, verifica los conteos y comprueba por patrones la
presencia de cada construcción de complejidad. Así se evita que la documentación y
la demostración se separen del comportamiento actual.

\section{Pruebas de distribución y recuperación de errores}

Además de los tres casos integrados en la interfaz, el directorio
\path{examples/error_distribution} contiene entradas manuales destinadas a
observar propagación, recuperación y control de cascadas. Estos archivos funcionan
como evidencia reproducible complementaria a las 24 pruebas Vitest para la
interfaz gráfica, la CLI y el video.

\begin{table}[H]
\centering
\caption{Resultados verificados de los casos de distribución de errores.}
\label{tab:distribucion-errores}
\begin{tabularx}{\textwidth}{>{\ttfamily}X r r r r}
\toprule
\normalfont Archivo & Tokens & Léxicos & Sintácticos & Total \\
\midrule
lexer\_errors.cps & 39 & 3 & 0 & 3 \\
parser\_errors.cps & 64 & 0 & 7 & 7 \\
lexer\_parser\_alternating\_errors.cps & 72 & 4 & 4 & 8 \\
\bottomrule
\end{tabularx}
\end{table}

\subsection{Caso exclusivamente léxico}

\archivo{lexer\_errors.cps} contiene los símbolos inválidos \archivo{@},
\archivo{\#} y \archivo{\$}. El fragmento \archivo{\#\#\#} se agrupa en un solo
diagnóstico. El resultado final contiene tres errores léxicos, ningún error
sintáctico derivado y tokens reconocidos después de cada fragmento incorrecto.

\subsection{Caso exclusivamente sintáctico}

\archivo{parser\_errors.cps} utiliza únicamente caracteres reconocibles, pero omite
puntos y coma y un paréntesis de cierre en construcciones independientes. El lexer no
reporta problemas y el parser produce siete diagnósticos en español, alcanza EOF y
conserva un árbol parcial mediante \archivo{DefaultErrorStrategy}.

\subsection{Caso léxico y sintáctico alternado}

\archivo{lexer\_parser\_alternating\_errors.cps} intercala cuatro símbolos inválidos
---\archivo{@}, \archivo{\#}, \archivo{\$} y un \archivo{\&} aislado--- con cuatro
omisiones sintácticas deliberadas. El resultado contiene exactamente cuatro errores
de cada fase. Los errores sintácticos inmediatos que podrían derivarse de los símbolos
descartados no se duplican, y el análisis continúa hasta la sentencia final
\archivo{print(e);}.

Los tres casos pueden ejecutarse desde la raíz del proyecto:

{\small
\begin{verbatim}
npm run cli -- examples/error_distribution/lexer_errors.cps
npm run cli -- examples/error_distribution/parser_errors.cps
npm run cli -- examples/error_distribution/lexer_parser_alternating_errors.cps
\end{verbatim}
}

Todos terminan con código de salida 1 porque son archivos intencionalmente inválidos;
ese código representa un rechazo esperado y no una excepción del analizador.

\section{Comandos de verificación}

La secuencia reproducible del proyecto es:

\begin{verbatim}
npm install
npm run generate
npm run check
npm test
npm run build
\end{verbatim}

\archivo{npm run check} ejecuta TypeScript sin emitir archivos. \archivo{npm test}
ejecuta Vitest una vez y \archivo{npm run build} combina comprobación de tipos y
construcción optimizada de Vite. Estas tres verificaciones cubren tipos, comportamiento
y empaquetado web, respectivamente.

% ================================================================
\chapter{Instalación, ejecución y distribución}
% ================================================================

\section{Requisitos de desarrollo}

\begin{itemize}
  \item Node.js 20.19 o posterior, o 22.12 o posterior;
  \item npm 9 o posterior;
  \item Java/JRE 11 o posterior únicamente para regenerar los analizadores;
  \item navegador moderno; Windows x64 para los ejecutables distribuidos.
\end{itemize}

\section{Ejecución web}

En la raíz del repositorio:

\begin{verbatim}
npm install
npm run generate
npm run dev -- --force
\end{verbatim}

Luego se abre \url{http://localhost:3000}. En ejecuciones posteriores normalmente
basta \archivo{npm run dev}. La opción \archivo{--force} resulta útil si Vite debe
reconstruir su caché después de instalar o cambiar dependencias.

\vspace{8cm}

\section{Uso de la aplicación}

\begin{enumerate}
  \item Seleccionar \textbf{Análisis completo} o \textbf{Solo lexer}.
  \item Elegir un caso incluido, escribir código o cargar un archivo \cps.
  \item Presionar \textbf{Analizar con ANTLR}.
  \item Revisar estado, tokens, ubicación de errores y, en modo completo, árbol.
  \item Descargar la entrada o los resultados si se necesita evidencia externa.
\end{enumerate}

\section{Uso de la CLI}

\begin{verbatim}
npm run cli:valid
npm run cli:error-lex
npm run cli:error-syn
npx tsx src/cli/run.ts ruta/al/programa.cps
\end{verbatim}

Los scripts de casos erróneos terminan intencionalmente con código 1; esto indica que
la herramienta rechazó correctamente la entrada, no que el analizador falló.

\section{Construcción de ejecutables}

\begin{verbatim}
npm run exe:portable
npm run exe:installer
\end{verbatim}

Ambos comandos construyen primero la aplicación web y luego invocan
\archivo{electron-builder}. Los artefactos quedan en \archivo{release}. Para una
entrega normal se publica el portable o el instalador. Los metadatos
\archivo{latest.yml}, los bloques \archivo{.blockmap}, los archivos internos NSIS y
la carpeta \archivo{win-unpacked} no forman parte de la entrega para usuarios.

Windows SmartScreen puede advertir sobre el ejecutable porque el proyecto académico
no posee un certificado comercial de firma. La procedencia debe comprobarse desde la
sección oficial de Releases antes de elegir ``Más información'' y ``Ejecutar de todas
formas''.

% ================================================================
\chapter{Cumplimiento del laboratorio}
% ================================================================

{\small
\begin{longtable}{p{0.32\textwidth} p{0.59\textwidth}}
\caption{Trazabilidad entre requisitos e implementación.}\label{tab:cumplimiento}\\
\toprule
Requisito & Evidencia en el proyecto \\
\midrule
\endfirsthead
\toprule
Requisito & Evidencia en el proyecto \\
\midrule
\endhead
Analizar Compiscript & Gramática \archivo{src/grammars/Compiscript.g4} y clases generadas. \\
Usar un generador & ANTLR 4 mediante \archivo{antlr4ts-cli}. \\
Recibir archivos \cps & Selector y validación de extensión en \archivo{InputEditor.tsx}; validación equivalente en CLI. \\
Integrar lexer y parser & \archivo{analyzeInput()} ejecuta ambos en el modo principal. \\
Mostrar resultados & Resumen, tabla de tokens, grupos de errores y árbol textual/visual. \\
Clasificar errores & Listas independientes \archivo{lexicalErrors} y \archivo{syntaxErrors}. \\
Indicar ubicación & Listener con línea y columna base 1. \\
Indicar símbolo & Campo \archivo{offendingSymbol} visible y serializable. \\
Mensajes inteligibles & Traducción y especialización en \archivo{antlrErrors.ts}. \\
Continuar tras un error & Tokenización completa y \archivo{DefaultErrorStrategy}. \\
Evitar repeticiones & Deduplicación, agrupación, filtrado de cascadas y límite por fase. \\
Confirmar éxito & Banner de aceptación cuando ambas listas están vacías. \\
Interfaz amigable & React, componentes especializados y diseño adaptable. \\
Limitar el alcance & Sin semántica, ejecución ni generación de código. \\
Evidencia práctica & Ocho casos oficiales de la rúbrica, tres casos de distribución, tres casos integrados y 24 pruebas Vitest. \\
\bottomrule
\end{longtable}
}

La trazabilidad muestra que los requisitos funcionales del laboratorio poseen una
implementación identificable y evidencia automatizada o visible. La UI no sustituye
al generador: presenta resultados provenientes de las clases creadas por ANTLR.

% ================================================================
\chapter{Resultados, limitaciones y trabajo futuro}
% ================================================================

\section{Resultados obtenidos}

El proyecto entrega un frente de compilador reutilizable en tres entornos. La
gramática reconoce las construcciones previstas de Compiscript, el lexer conserva
tokens posteriores a símbolos inválidos y el parser produce un árbol incluso después
de recuperaciones posibles. La normalización hace que la salida sea más apropiada
para una persona que aprende el lenguaje, sin ocultar la fase ni la ubicación del
problema.

La separación entre motor, tipos y presentación facilita las pruebas y evita que las
reglas de aceptación dependan de React o Electron. Los ejemplos sincronizados reducen
el riesgo de que la demostración visual difiera de los archivos entregados.

\section{Limitaciones actuales}

\begin{itemize}
  \item La corrección se define solo a nivel léxico y sintáctico.
  \item La recuperación puede producir un árbol parcial con tokens insertados o
        descartados conceptualmente; ese árbol no debe confundirse con ejecución.
  \item El límite de 100 diagnósticos prioriza una salida manejable sobre mostrar cada
        consecuencia de una entrada extremadamente dañada.
  \item El filtrado de cascadas usa cercanía de posición; es una mejora de
        presentación, no un análisis causal formal de errores.
  \item Los ejecutables configurados se dirigen a Windows x64.
  \item La firma comercial de código no forma parte del proyecto académico.
\end{itemize}

\section{Trabajo futuro}

Una evolución natural agregaría una tabla de símbolos y un visitor para comprobar
declaraciones, ámbitos, herencia, llamadas y tipos. Después podría incorporarse un
intérprete o una representación intermedia. También sería posible enriquecer los
diagnósticos con fragmentos de fuente y subrayado, pruebas de propiedades para la
gramática, medición de cobertura y empaquetado multiplataforma firmado.

% ================================================================
\chapter{Conclusiones}
% ================================================================

\begin{enumerate}
  \item Una gramática declarativa permitió generar un lexer y parser reales y
        mantener separadas la definición del lenguaje y la interfaz de usuario.
  \item Forzar la tokenización completa antes del parseo es esencial para recolectar
        múltiples errores léxicos y demostrar que el análisis continúa.
  \item La recuperación predeterminada de ANTLR, acompañada de deduplicación,
        agrupación y filtrado, produce diagnósticos más útiles que detenerse en la
        primera anomalía o mostrar cascadas sin control.
  \item La jerarquía de reglas de expresión implementa precedencia sin construir un
        parser manual ni usar tablas SLR.
  \item Un contrato único de resultado hizo posible reutilizar el analizador desde
        React, la CLI y Electron, además de facilitar exportaciones y pruebas.
  \item Las 24 pruebas automatizadas y la matriz de ocho casos de la rúbrica aportan
        evidencia de cobertura, recuperación, continuidad y consistencia.
  \item El proyecto cumple el propósito del Laboratorio 01 dentro de su alcance
        léxico y sintáctico; las etapas semánticas y de generación quedan claramente
        separadas como trabajo posterior.
\end{enumerate}

\begingroup
\renewcommand{\bibname}{Referencias}
\addcontentsline{toc}{chapter}{Referencias}
\small
\setlength{\parskip}{0pt}
\begin{thebibliography}{9}
\setlength{\itemsep}{0.45em}
\setlength{\parsep}{0pt}

\bibitem{aho}
Alfred V. Aho, Monica S. Lam, Ravi Sethi y Jeffrey D. Ullman.
\textit{Compilers: Principles, Techniques, and Tools}.
Addison-Wesley, segunda edición, 2006.

\bibitem{parr}
Terence Parr.
\textit{The Definitive ANTLR 4 Reference}.
Pragmatic Bookshelf, segunda edición, 2013.

\bibitem{antlr}
ANTLR Project.
\textit{ANTLR 4 Documentation}.
\url{https://www.antlr.org/}

\bibitem{typescript}
Microsoft.
\textit{TypeScript Documentation}.
\url{https://www.typescriptlang.org/docs/}

\bibitem{react}
React Team.
\textit{React Documentation}.
\url{https://react.dev/}

\bibitem{electron}
OpenJS Foundation.
\textit{Electron Documentation}.
\url{https://www.electronjs.org/docs/}

\bibitem{repo}
DragonsSlayers 2.0.
\textit{Analizador de Compiscript: código fuente, README, gramática, pruebas y documentación técnica}.
Repositorio del Laboratorio 01, 2026.

\end{thebibliography}
\endgroup

% ================================================================
\appendix
\chapter{Preguntas interesantes a las cuales se responde}
% ================================================================

\begin{description}
  \item[¿El analizador es real?] Sí. La función principal instancia el lexer y el
  parser de Compiscript generados por ANTLR e invoca \regla{program}.

  \item[¿Cómo continúa después de errores?] El lexer completa el flujo mediante
  \archivo{fill()} y el parser usa \archivo{DefaultErrorStrategy}; los listeners
  acumulan diagnósticos en lugar de lanzar la interfaz al primer error.

  \item[¿Cómo se evita repetir mensajes?] Se deduplican coincidencias exactas, se
  agrupan caracteres contiguos, se omiten consecuencias sintácticas inmediatas de
  errores léxicos y se limita la vista por fase.

  \item[¿Por qué no hay tabla SLR?] ANTLR 4 genera un parser Adaptive LL(*), por lo
  que no utiliza el algoritmo SLR ni sus tablas \archivo{ACTION/GOTO}.

  \item[¿Qué prueba la aceptación?] Que no hubo errores léxicos ni sintácticos. No
  prueba corrección de tipos ni comportamiento en ejecución.

  \item[¿Por qué existe el código generado?] Es el resultado reproducible de aplicar
  ANTLR a la gramática. Se versiona para ejecutar el proyecto sin exigir Java al
  usuario final, pero se regenera después de cambiar \archivo{Compiscript.g4}.
\end{description}

\chapter{Inventario de comandos}

\begin{tabularx}{\textwidth}{>{\ttfamily}l X}
\toprule
Comando & Propósito \\
\midrule
npm run generate & Regenerar lexer, parser, listener y visitor. \\
npm run dev & Iniciar Vite en el puerto 3000. \\
npm run check & Verificar tipos TypeScript sin emitir archivos. \\
npm test & Ejecutar las 24 pruebas una vez. \\
npm run build & Validar tipos y crear \archivo{dist}. \\
npm run cli:valid & Analizar el ejemplo correcto desde terminal. \\
npm run cli:error-lex & Demostrar errores léxicos; salida esperada 1. \\
npm run cli:error-syn & Demostrar recuperación sintáctica; salida esperada 1. \\
npm run cli -- archivo.cps & Analizar cualquier caso manual de la carpeta de ejemplos. \\
npm run desktop & Construir web y abrir Electron. \\
npm run exe:portable & Crear ejecutable portable Windows x64. \\
npm run exe:installer & Crear instalador NSIS Windows x64. \\
\bottomrule
\end{tabularx}

\chapter{Ficha técnica del producto}

\begin{tabularx}{\textwidth}{>{\bfseries}l X}
\toprule
Elemento & Valor \\
\midrule
Lenguaje analizado & Compiscript \\
Extensión de entrada & \cps \\
Gramática & \archivo{src/grammars/Compiscript.g4} \\
Generador & ANTLR 4 mediante \archivo{antlr4ts-cli} \\
Lenguaje objetivo & TypeScript \\
Frontend & React 19 y Vite 6 \\
Pruebas & Vitest 3.2, 24 casos \\
Escritorio & Electron 43 y electron-builder 26 \\
Plataforma empaquetada & Windows x64 \\
Motor visible & ANTLR 4 / antlr4ts, código generado \\
Regla inicial & \regla{program} \\
Recuperación & \archivo{DefaultErrorStrategy} \\
Máximo visible & 100 diagnósticos por fase \\
Salidas & UI, consola, CSV, JSON, TXT y árbol visual \\
Alcance & Análisis léxico y sintáctico \\
\bottomrule
\end{tabularx}

\vfill
\begin{center}
\color{uvgblue}\rule{0.75\textwidth}{0.8pt}\color{black}

\textbf{DragonsSlayers 2.0}\par
Laboratorio 01 -- Construcción de Compiladores CC-3032\par
Documentación elaborada a partir de la versión actual del repositorio.
\end{center}

\end{document}
